\section{Localization to one transversal orientation}
\label{app:code-orientation-localization}

This appendix treats arbitrary coordinate weights, without assuming
that they are constant on a rail. There is one shared carrier coordinate
and $m\ge2$ disjoint pairs of even rails. Every rail size belongs to
$[R,R+\Delta]$, where $R\ge4$ is even and $0\le\Delta\le R$.
A full transversal support consists of the carrier and one rail from
every pair. Any retained subfamily $\mathcal S$ is allowed. Put $D=mR$
and
\[
 Q(h_u)=\sum_{s\in\mathcal S}w_{n_s}|\widehat h_u(A_s)|,
 \qquad h_u(x)=\tanh(u\cdot\bar x),\qquad n_s=|A_s|.
\]
The incoming norm includes the bias and any irrelevant coordinates.
The balanced scalar estimates give $w_n\le Cn$ for odd $n\ge3$;
hence $w_{\max}\le CD$ here.

\begin{theorem}[A single retained orientation carries the mass]
\label{thm:code-orientation-localization}
Fix $\beta>0$. There are positive constants $c,C,K_0$, depending only
on $\beta$, with the following property. Suppose $u\ne0$,
$R\ge K_0\log^2(m+2)$, and $Q(h_u)\ge\beta\|u\|_2$.
Then
\[
 c\sqrt D\le\|u\|_2\le C\sqrt D,
\]
and there is a retained orientation $s_*$ such that
\begin{equation}\label{eq:code-off-orientation}
 \sum_{s\ne s_*}|\widehat h_u(A_s)|
 \le Cm e^{-cR}
   +\frac{C(1+\sqrt R)e^{-c\sqrt R}}{\sqrt D}.
\end{equation}
The sum may run over the whole full code, so the same bound holds on
the retained subfamily. For every fixed $L>0$, there is a constant
$K_L\ge K_0$ such that all sufficiently large $m$ satisfy, uniformly
in $R\ge K_L\log^2(m+2)$,
\begin{equation}\label{eq:code-off-superpolynomial}
 \sum_{\substack{s\in\mathcal S\\s\ne s_*}}
 w_{n_s}|\widehat h_u(A_s)|\le D^{-L}\|u\|_2.
\end{equation}
Both $K_L$ and the threshold for $m$ may depend on $L$ and $\beta$.
Only the fixed choice $L=8$ is needed in the exact code-norm theorem.
\end{theorem}

\begin{proof}
Input sign changes allow assuming nonnegative coordinate magnitudes.
For pair $i$ define
\[
 A_i(t)=\prod_{j\in B_i^0}|\sin(u_jt)|
        \prod_{j\in B_i^1}|\cos(u_jt)|,\qquad
 B_i(t)=\prod_{j\in B_i^0}|\cos(u_jt)|
        \prod_{j\in B_i^1}|\sin(u_jt)|,
\]
and put $F(t)=\prod_i(A_i(t)+B_i(t))$.
Write $F_{\mathcal S}$ for the corresponding sum of products over
retained orientations. The coefficient integral for an odd support
contains a common carrier sine, a bias cosine, and irrelevant-coordinate
cosines. Taking absolute values therefore gives
\begin{equation}\label{eq:code-positive-integral}
 \sum_{s\in\mathcal S}|\widehat h_u(A_s)|
 \le\int_0^\infty
       \frac{|\sin(u_0t)|F_{\mathcal S}(t)}{\sinh(\pi t/2)}\,dt.
\end{equation}
Here $u_0$ is the carrier magnitude. Each $A_i+B_i\le1$: keep one
coordinate from each rail, discard the other factors, and apply
$|\sin x\cos y|+|\cos x\sin y|\le1$. Thus $0\le F\le1$.

\paragraph{Uniform scale bounds.}
Choose any $R$ coordinates from each rail and let $L_0$ be their combined
Euclidean norm. Put $\lambda=L_0/\sqrt D$. We claim
\begin{equation}\label{eq:code-scale-integral}
 \int_0^\infty\frac{F(t)}{\sinh(\pi t/2)}\,dt
 \le\frac C{\sqrt D}
       \min\{\lambda^3,\log(1+\lambda)\}.
\end{equation}
Discard rail factors beyond the chosen coordinates and zip each pair
into $R$ columns. Generalized H\"older over the orientation index gives
an upper bound by the geometric mean of the column functions
\[
 F_k(t)=\prod_{i=1}^m\left[
 \sin(a_it)^R\cos(b_it)^R+\cos(a_it)^R\sin(b_it)^R\right].
\]
The notation $a_i,b_i$ refers to a fixed column; the two magnitudes may
be ordered separately in each column. Clarkson's even-power inequality
bounds this function by
\[
 \E_{\varepsilon}\prod_i|\sin(c_it)|^R,
 \qquad c_i=|a_i+\varepsilon_i b_i|.
\]
H\"older over the $m$ pairs bounds its integral by
$\E_\varepsilon\prod_i r_D(c_i)^{1/m}$, where $r_D$ is the absolute
diagonal integral used in Appendix~\ref{app:odd-lattice-localization}.
Equation~\eqref{eq:absolute-diagonal-comparison} and the elementary
bounds on $r_3$ give
\[
 r_D(c)\le\frac C{\sqrt D}r_3(c)
 \le\frac C{\sqrt D}\min\{c^3,\log(1+c)\}.
\]
For the logarithmic bound, geometric mean at most arithmetic mean and
concavity of $s\mapsto\log(1+\sqrt s)$ imply
\[
 \E_\varepsilon\prod_i\log(1+c_i)^{1/m}
 \le\log\left(1+\sqrt{\frac1m\sum_i(a_i^2+b_i^2)}\right).
\]
For the cubic bound use geometric mean at most root mean square and
$\sum_i c_i^2\le2\sum_i(a_i^2+b_i^2)$ pointwise. Finally apply
H\"older over the $R$ columns. The geometric mean of the cubic bounds
is at most the cube of the combined root mean square; for the logarithmic
bounds use the same concavity once more. This proves
\eqref{eq:code-scale-integral}.

Set $z=\|u\|_2/\sqrt D$. Since $L_0\le\|u\|_2$, the functions in
the right side of \eqref{eq:code-scale-integral} are increasing, and
$w_{\max}\le CD$. The competitiveness assumption yields
\[
 \beta z\le C\min\{z^3,\log(1+z)\}.
\]
Consequently $z$ belongs to a fixed compact interval
$[z_-,z_+]\subset(0,\infty)$, proving the asserted norm bounds.

\paragraph{A constant-height resonance.}
For every $T\ge0$, every frequency $v\ge0$, and every $p\ge1$,
\begin{equation}\label{eq:uniform-sine-window}
 \int_T^{T+1}|\sin(vt)|^p\,dt\le\frac{C(T+1)}{\sqrt p}.
\end{equation}
If $v\le\pi/[4(T+1)]$, the sine is at most $1/\sqrt2$ there.
Otherwise substitution and covering by full or partial periods proves
the estimate, using $1/v\le4(T+1)/\pi$.
The column argument above also gives
\begin{equation}\label{eq:code-half-window}
 \int_T^{T+1}F(t)^{1/2}\,dt\le\frac{C(T+1)}{\sqrt D}.
\end{equation}
Indeed H\"older across columns reduces to a single column. H\"older
across the $m$ pairs then produces powers $m/2$ of a single pair factor.
Since $m/2\ge1$, Clarkson followed by Jensen bounds these by averages
of $|\sin((a_i\pm b_i)t)|^{D/2}$, to which
\eqref{eq:uniform-sine-window} applies. As $F\le1$, the same bound
holds for $F$ itself.

It follows that the weighted integral beyond $T\ge1$ is at most
$C(T+1)e^{-\pi T/2}/\sqrt D$. Near zero, every transversal uses at
least $D$ rail coordinates, and arithmetic--geometric mean gives
\[
 F(t)\le2^m(z_+t)^D\qquad(z_+t\le1).
\]
A fixed small lower cutoff therefore makes the near-zero contribution
exponentially small in $D$. On the other hand,
\eqref{eq:code-positive-integral}, competitiveness, and $w_{\max}\le CD$
show that the retained positive integral is at least $c/\sqrt D$.
Choose fixed $0<t_-<t_+<\infty$ so that
\[
 \int_{t_-}^{t_+}
 \frac{|\sin(u_0t)|F_{\mathcal S}(t)}{\sinh(\pi t/2)}\,dt
 \ge\frac c{\sqrt D}.
\]
Put $P(t)=|\sin(u_0t)|F_{\mathcal S}(t)$ and
$M=\max_{[t_-,t_+]}P$. Since $P\le F$ and $P\le M$, one has
$P\le\sqrt M\sqrt F$. Equation~\eqref{eq:code-half-window}, with
the bounded kernel on this compact interval, therefore gives
$M\ge\mu>0$, independently of $m,R$. Choose $t_*$ attaining this bound.

\paragraph{One retained orientation.}
For each pair,
\[
 A_i(t_*)B_i(t_*)\le2^{-(r_i^0+r_i^1)}\le4^{-R},
 \qquad \min\{A_i(t_*),B_i(t_*)\}\le2^{-R}.
\]
Select the larger factor in every pair and call the resulting orientation
$s_*$. Since all pair sums are at most one, telescoping products gives
\[
 P_{s_*}(t_*):=\prod_i\max\{A_i(t_*),B_i(t_*)\}
 \ge F(t_*)-m2^{-R}\ge\mu-m2^{-R}.
\]
For $R\ge K_0\log^2(m+2)$ and sufficiently large $m$, this is at least $\mu/2$.
The orientation is retained: otherwise
$F_{\mathcal S}(t_*)\le F(t_*)-P_{s_*}(t_*)\le m2^{-R}<\mu$,
contrary to $P(t_*)\ge\mu$.

Let $a=\pi/(2t_*)$, in a fixed compact positive interval. Choose nearest
integers $k_j$ of odd parity on selected rails and even parity on
unselected rails, and write
\[
 \frac{2u_jt_*}{\pi}=k_j+\delta_j.
\]
Taking minus logarithms of $P_{s_*}(t_*)\ge\mu/2$, and using
$-\log\cos x\ge x^2/2$, shows that
\begin{equation}\label{eq:code-old-phase-energy}
 \sum_{\text{rail }j}\delta_j^2\le K
\end{equation}
for a fixed constant $K$. Thus the selected coordinates are close in
total squared distance to an odd lattice, and the unselected ones to
its even lattice.

\paragraph{Whole-set parity determinants.}
The norm bound and \eqref{eq:code-old-phase-energy} also give
\[
 \sum_{\text{rail }i}k_i^2\le K_1D
\]
for a fixed constant $K_1$. We use this total energy bound instead of
bounding individual coordinates separately in each rail.
Fix another orientation which switches exactly $j$ pairs. At time
$t>0$, put $\gamma=t/t_*$ and choose nearest integers of the new
required parities, writing
\[
 \gamma(k_i+\delta_i)=l_i+\epsilon_i,\qquad
 E=\sum_{\text{rail }i}\epsilon_i^2.
\]
For disjoint sets of coordinates $A,B$ such that every determinant
$k_a l_b-k_b l_a$ is odd, substitution and squaring give
\begin{align*}
 |A||B|&\le\sum_{a\in A,b\in B}(k_a l_b-k_b l_a)^2\\
 &\le4\left[K_AE_B+K_BE_A+
       \gamma^2(K_A\Delta_B+K_B\Delta_A)\right]\\
 &\le4K_1D(E+K\gamma^2),
\end{align*}
where $K_A=\sum_Ak_i^2$, $E_A=\sum_A\epsilon_i^2$, and
$\Delta_A=\sum_A\delta_i^2$.
Indeed the determinant equals
$\gamma(k_a\delta_b-k_b\delta_a)-k_a\epsilon_b+k_b\epsilon_a$.

If $j\le m/2$, take $A$ to be the old selected coordinates in switched
pairs and $B$ the old selected coordinates in unchanged pairs.
Their old parities are both odd, while their new parities are even
and odd, respectively. Thus every determinant is odd, and
$|A|\ge jR$, $|B|\ge(m-j)R$.
If $j>m/2$, take $A$ as before and $B$ the old unselected coordinates
in switched pairs. Their parity types are odd-to-even and even-to-odd;
again every determinant is odd, and both sets have size at least $jR$.
In both cases
\[
 E\ge c jR-C\gamma^2.
\]
Since $t_*$ is bounded away from zero, there is a fixed $c_1>0$ such
that $E\ge c_2jR$ whenever $0<t\le c_1\sqrt R$.
Each sine or cosine factor is $|\cos(\pi\epsilon_i/2)|$, with
$|\epsilon_i|\le1$, and this is at most
$e^{-\pi^2\epsilon_i^2/8}$. Hence the product of an orientation
with $j$ switches is at most $e^{-c_3jR}$ throughout this time range.
Summing by the number of switches gives
\[
 F(t)-P_{s_*}(t)
 \le\sum_{j=1}^m\binom mj e^{-c_3jR}
 =(1+e^{-c_3R})^m-1\le Cm e^{-c_3R},
\]
where the final inequality holds for $R\ge K_0\log^2(m+2)$ and
large $m$.

Between a fixed lower cutoff and $c_1\sqrt R$, the Fourier kernel
has bounded total integral, so this contributes $Cm e^{-cR}$.
Below the cutoff, the earlier arithmetic--geometric-mean bound is
exponentially small in $D$ and is absorbed by this term. Above
$c_1\sqrt R$, use the full-code tail from
\eqref{eq:code-half-window}. This proves
\eqref{eq:code-off-orientation}.

Multiplying by $w_{\max}\le CD$ and dividing by
$\|u\|_2\ge c\sqrt D$ bounds the relative off-mass by
\[
 C\sqrt D\,m e^{-cR}+C(1+\sqrt R)e^{-c\sqrt R}.
\]
Here $D=mR$. If $R\ge K_L\log^2(m+2)$, then
$\log m\le\sqrt{R/K_L}$, while $\log R=o(\sqrt R)$ uniformly
as $m\to\infty$. Choose the fixed $K_L$ large enough compared with
$L^2/c^2$. The last exponential then absorbs every factor in
$D^L(1+\sqrt R)$; the first term is smaller still. This proves
\eqref{eq:code-off-superpolynomial}, with the stated uniformity.
\end{proof}
